\section{Stationary equilibrium}\label{sec:eqm}

Every cohort is born with no assets and an earnings state drawn from a fixed law $\nu$, lives $J$
periods, and is replaced; cohorts have equal weight. The distribution of the youngest cohort is
therefore the same at every interest rate, and the distribution of every older cohort is that law
pushed forward by the policies of the ages in between. This is the structural advantage of the
life cycle over the infinite horizon, and it removes most of the machinery the stationary case
needed.

Two of the consumption bounds this section uses are developed in Section~\ref{sec:thm1}, where
they arise in the hunt for the critical risk aversion: the constraint-incidence bound of
Theorem~\ref{thm:incidence} and the asset-dependent saving floor of Section~\ref{sec:assetdep}.
Nothing here depends on the verdict of that section.

\subsection{The aggregate is a finite composition}

Because assets at birth are degenerate at zero, no measure theory is required. Let
\[
  (\mathcal{S}_k h)(a,z)=\sum_{z'}\pi(z,z')\,h\bigl(g_k(a,z),z'\bigr)
\]
be the expectation of $h$ next period for a household with $k$ periods still to come
(\lean{stageStep}). The assets a household holds over the rest of its life, summed across ages
and seen from its current state, satisfy
\[
  A_0(a,z)=a,\qquad A_{k+1}(a,z)=a+(\mathcal{S}_{k+1}A_k)(a,z)
\]
(\lean{cohortAssets}), and aggregate capital per head is
$K=\frac1J\sum_z\nu(z)\,A_{J-1}(0,z)$ (\lean{olgCapital}). There is no pushforward measure, no
Feller property, no stationary distribution, no Doeblin condition and no uniqueness of an
invariant law: the equilibrium object is a finite sum of finite sums, and the recursion above is
its definition.

\subsection{The chain}

\begin{theorem}[Light's Theorem 2 for the life cycle]\label{thm:thm2}
Suppose two economies share the earnings chain and their policies are ordered at every age,
$g^{(1)}_k\le g^{(2)}_k$ on the asset region. Then $A^{(1)}_k\le A^{(2)}_k$ for every $k$, and
aggregate capital is ordered, $K^{(1)}\le K^{(2)}$.
\formal{IncomeFluctuation.cohortAssets\_le,
IncomeFluctuation.olgCapital\_le\_of\_stagePolicy\_le}
\end{theorem}

The proof is an induction on age with two ingredients, and both are already available. The class
of functions increasing in assets is preserved by $\mathcal{S}_k$
(\lean{monoRegion\_stageStep}), because next period's assets rise with this period's at each
earnings state and for any continuation, which is Theorem~\ref{thm:azmono}. And ordered policies
keep ordered functions ordered (\lean{stageStep\_le\_stageStep}): the higher-rate economy
evaluates a larger function at a larger asset level. Note what is \emph{not} needed. The
infinite-horizon version of this step required monotone transitions, a coupling argument and
first-order stochastic dominance of the stationary distributions; here the whole statement is an
inequality between two functions on the same finite state space, because the distribution the two
economies start from is literally the same.

\subsection{Uniqueness}

\begin{theorem}[Single crossing]\label{thm:thm3}
A capital supply that never falls with the rate meets a strictly falling demand at most once.
\formal{IncomeFluctuation.olgEquilibriumRate\_unique}
\end{theorem}

With Theorem~\ref{thm:thm2} this reduces uniqueness of the stationary equilibrium to Light's
Theorem 1, that saving rises with the rate at every age, which is the subject of
Section~\ref{sec:thm1}: it holds at Aiyagari's calibration for relative risk aversion up to about
$3.4$, so the chain closes at his $\mu=3$ and not at his $\mu=5$.

\subsection{Continuity, and existence}

The last input is continuity of capital supply in the rate. In the infinite horizon this was a
substantial development, because the stationary distribution itself had to be shown to move
continuously. Here there is no stationary distribution: capital supply is a finite composition of
finitely many policies, so all that is needed is that each policy moves continuously with the
rate, and that continuity composes.

We take the route of the companion write-up and put the interest rate into the state, so that
Berge's theorem gives continuity in assets and the rate jointly rather than continuity in a
parameter. That device was used there at the fixed point, but everything it rests on --- that the
feasible sets and the objectives agree on the slice --- holds for an arbitrary continuation. So
it works stage by stage, and here it is easier, because the slice of the augmented Bellman
iterate is the iterate of the sliced programme by a plain induction, with no fixed point to
identify (\lean{sliceAt\_bellman}, \lean{sliceAt\_augStageValue}). The augmented argmax at each
stage is then a singleton, since the sliced one is by Carroll--Kimball concavity of the iterates
(\lean{argmax\_aug\_stage\_eq\_singleton}), so the stage policy is jointly continuous
(\lean{continuousOn\_augStagePolicy}) and agrees with the sliced one
(\lean{augStagePolicy\_eq}). Running the cohort recursion in the augmented state
(\lean{augCohortAssets}) gives a composition of continuous maps that reads, at a rate in range, as
capital supply at that rate (\lean{augCohortAssets\_eq}, \lean{olgSupply\_eq}).

\begin{theorem}[Capital supply is continuous]\label{thm:cont}
$r\mapsto K(r)$ is continuous on $[r_{lo},r_{hi}]$.
\formal{IncomeFluctuation.continuousOn\_olgSupply}
\end{theorem}

\begin{theorem}[Existence]\label{thm:exist}
If capital supply lies below demand at the bottom of the interval and above it at the top, and
demand is continuous, then the two meet: there is a rate in $[r_{lo},r_{hi}]$ at which the capital
the cohorts accumulate equals the capital the firm demands.
\formal{IncomeFluctuation.exists\_olgEquilibrium}
\end{theorem}

Both hypotheses of Theorem~\ref{thm:exist} are about the solved model. The rest of this section
replaces them with inequalities that never mention the household's policy or value function.

\subsection{Where the household can be}\label{sec:reach}

Proposition~\ref{prop:floor} inherits from Theorem~\ref{thm:sandwich} the hypothesis that the
artefactual cap is slack, and the convenient way to state that is slack everywhere in
$[0,\bar a]$. That hypothesis is not innocent, and it fails exactly where the floor is wanted.

The one-step bound gives $g_k(a,z)\le(1-\kappa_k)(Ra+y_{\max})$, so $[0,\bar a]$ is carried into
itself only if $(1-\kappa_k)(R\bar a+y_{\max})\le\bar a$, which needs $(1-\kappa_k)R<1$. As
$k$ grows $\kappa_k$ falls to $1-\text{\th}/R$, so the requirement is $\text{\th}<1$, that is
$\beta R<1$: the household must be impatient. Above the rate of time preference no cap is forward
invariant at any level, because a patient household sitting at the cap wants to save past it. The
cap then binds, the Euler inequality under it is unavailable, and the slackness hypothesis is
false rather than merely unverified. A theorem carrying it says nothing there.

This matters because a stationary life-cycle equilibrium sits above the rate of time preference.
In the calibration of Section~\ref{sec:num2} with $\lambda=1/\beta-1=4.17\%$, capital supply
crosses demand near $5\%$, and the rate at which the \emph{provable} floor crosses demand is
near $9\%$. Both are in the region where no cap is slack everywhere.

A cohort that begins life with nothing never visits the top of $[0,\bar a]$. The development
therefore carries a family of regions, one per stage, that the stage policies map down through,
and asks for slackness only on those.

\begin{definition}[Reachable family]\label{def:regions}
A reachable family assigns to each stage $k$ a set of asset levels such that every set sits
inside $[0,\bar a]$, the stage-$(k+1)$ policy carries the stage-$(k+1)$ set into the stage-$k$
set, and $g_k<\bar a$ on every set.
\formal{IncomeFluctuation.StageRegions}
\end{definition}

The whole interval $[0,\bar a]$ is a reachable family exactly when the cap is slack on it
(\lean{allRegions}), so nothing is lost: Theorem~\ref{thm:sandwich} and
Proposition~\ref{prop:floor} are the constant-family instances of statements that hold on any
reachable family (\lean{stageConsumption\_le\_stageHumanWealth\_on}, \lean{le\_stagePolicy\_on}).

What a cohort visits is bounded by feasibility alone. Set $\rho_0=0$ and
$\rho_{j+1}=R\rho_j+y_{\max}$: a household that cannot carry forward more than its cash on hand
holds at most $\rho_j$ after $j$ periods of life, whatever the earnings draws.

\begin{proposition}[A cohort's own regions]\label{prop:reach}
If $\rho_{K+1}<\bar a$ then the sets $[0,\rho_{K-k}]$, empty for $k>K$, form a reachable family
for a cohort of length $K+1$ that starts life with nothing.
\formal{IncomeFluctuation.reachRegions}
\end{proposition}

The bound is deliberately crude, using none of the damping that optimal saving supplies, and that
is what keeps it free of the sandwich it is used to prove. Its cost is a larger cap. At $r=10\%$
over sixty periods $\rho_{K+1}$ is about $9{,}044$ where the optimal path reaches about $82$.
Since the cap is an artefact of the state space rather than an economic object, a cap of ten
thousand is as acceptable as a cap of one hundred, and unlike slackness everywhere it is a
condition that can actually be met. The sandwich and the floor then hold along a cohort's path
with no hypothesis on the cap beyond $\rho_{K+1}<\bar a$
(\lean{stageConsumption\_le\_stageHumanWealth\_reach}, \lean{le\_stagePolicy\_reach}).

At the bottom of the rate interval the same crude bound does the other half of the work on its
own. When $R<1$ the sequence $\rho_j$ converges to $y_{\max}/(1-R)$, so cohort assets are bounded
without any appeal to optimality: at $r=-7\%$ the bound is $42.1$ against a capital demand of
$56.3$.

\subsection{The propensity in closed form}\label{sec:geom}

Both bounds below are built from the finite-horizon propensity $\kappa_k$ and from human wealth,
and both need to know not just where $\kappa_k$ sits but how fast it settles.

\begin{proposition}[The finite-horizon propensity]\label{prop:kappa}
Write $q=\Thorn/R$. Then $\kappa_k\sum_{i\le k}q^i=1$, and hence
$\kappa_k\bigl(1-q^{k+1}\bigr)=1-q$: the propensity exceeds its infinite-horizon limit $1-q$ by
exactly $\kappa_kq^{k+1}$.
\formal{IncomeFluctuation.stageMPC\_mul\_mpcSum},
\lean{IncomeFluctuation.stageMPC\_mul\_one\_sub\_pow}
\end{proposition}

The defining recursion $\kappa_{k+1}=\kappa_kR/(\Thorn+\kappa_kR)$ is exactly
$1/\kappa_{k+1}=1+q/\kappa_k$, so the reciprocal accumulates the geometric series and nothing
else. The same sum tames the slope that both bounds share. Writing $S_k=\sum_{i\le k}q^i$, the
recursion $G_0=1$, $G_{k+1}=1+G_k(1-\kappa_{k+1})R$ has an age-dependent coefficient, but
$W_k=G_kS_k$ obeys $W_{k+1}=S_{k+1}+\Thorn W_k$, which does not. Then $S_{k+1}\ge1$ and
$(1-q)S_k\le1$ give the bound that matters.

\begin{proposition}[A slope that grows]\label{prop:slope}
$G_k\ \ge\ (1-q)\sum_{j\le k}\Thorn^{\,j}$.
\formal{IncomeFluctuation.mul\_geomSum\_le\_cohortSlope}
\end{proposition}

Where the crude $G_k\ge1$ does not grow at all, this grows geometrically, and
Section~\ref{sec:sharpnum} reports what that is worth.

\subsection{A floor on capital}\label{sec:clampfloor}

Existence needs capital supply above demand at the top of the rate interval. In the infinite
horizon the only floor available came from an ergodic identity that loses a factor
$\|u'\|_\infty^2/\operatorname{Var}$. The life cycle needs nothing of the kind: the upper sandwich
of Theorem~\ref{thm:sandwich} bounds consumption at every age, and what is not consumed is saved.

\begin{proposition}[Affine floor on saving]\label{prop:floor}
At every age and every state in the region,
\[
  g_k(a,z)\ \ge\ (1-\kappa_k)\bigl(y_z+Ra\bigr)-\kappa_k H_k(z).
\]
\formal{IncomeFluctuation.le\_stagePolicy\_on}, \lean{IncomeFluctuation.le\_stagePolicy}
\end{proposition}

Iterating it is not a matter of tracking a distribution over assets. An affine lower bound on a
cohort's accumulated assets suffices, $F_k(z)+G_ka\le A_k(a,z)$, with the slope $G$ of
Section~\ref{sec:geom} common to all earnings states and an intercept carrying one value per
state:
\[
  F_0=0,\qquad
  F_{k+1}(z)=\textstyle\sum_{z'}\pi(z,z')F_k(z')
    +G_k\bigl[(1-\kappa_{k+1})y_z-\kappa_{k+1}H_{k+1}(z)\bigr].
\]

\begin{proposition}[Affine floor on accumulated assets]\label{prop:affine}
On any reachable family, $F_k(z)+G_ka\le A_k(a,z)$; hence capital per head is at least
$\sum_z\nu_zF_K(z)/(K+1)$.
\formal{IncomeFluctuation.cohortFloor\_add\_mul\_le\_cohortAssets},
\lean{IncomeFluctuation.le\_olgCapital}
\end{proposition}

The state dependence of $F$ is the whole content, not a refinement. Reading
Proposition~\ref{prop:floor} at the worst earnings state instead, which is the only way to get a
bound uniform in $z$, gives a floor that is identically zero at every rate in the calibration of
Section~\ref{sec:num2}, against $4.75$ at $r=10\%$ for the state-dependent one.

Averaging against weights the chain leaves alone collapses the vector recursion to a scalar one,
and the sign of each increment then has an exact answer.

\begin{theorem}[The floor rises exactly when the household is patient]\label{thm:patient}
Let $\bar y$ and $\bar H_k$ be the averages of earnings and of human wealth under weights $\nu$
satisfying $\nu\pi=\nu$. If $\Thorn\ge1$, equivalently $\beta R\ge1$, then
$\kappa_k\bar H_k\le(1-\kappa_k)\bar y$ at every age, so the aggregate floor never falls and is
never negative.
\formal{IncomeFluctuation.stageMPC\_mul\_sum\_stageHumanWealth\_le},
\lean{IncomeFluctuation.sum\_mul\_cohortFloor\_mono}
\end{theorem}

The increment of the aggregate floor is $G_k\kappa_{k+1}\bar y$ times the slack in that
inequality, and the inequality reverses below $\Thorn=1$. So the provable floor is flat under the
rate of time preference and climbs above it, which is why the upper endpoint of the rate interval
has to be taken above $\lambda=1/\beta-1$. The proof is one induction: writing
$v_k=(1-\kappa_k)/\kappa_k$ and $u_k=\bar H_k/\bar y$, both satisfy first-order recursions,
$v_{k+1}=(\Thorn/R)(1+v_k)$ and $u_{k+1}=(1+u_k)/R$, from the same starting point, and
$\Thorn\ge1$ keeps the first above the second.

One thing the affine form cannot do is refuse to claim negative assets. Its intercept goes badly
negative: at the calibration, $160$ of $413$ entries are below zero and the worst is $-62$, with
$F_K$ running from $-58.8$ at the lowest earnings to $+685$ at the highest. A cohort's accumulated
assets cannot be negative, so all of that is slack the aggregate inherits. Reading the floor as a
function of assets rather than as an affine form fixes it, because the saving floor can then be
clamped at zero where it belongs:
\[
  \Lambda_0(z,a)=a,\qquad
  \Lambda_{k+1}(z,a)=a+\textstyle\sum_{z'}\pi(z,z')\,\Lambda_k\bigl(z',\Phi_{k+1}(z,a)\bigr),
\]
with $\Phi$ the clamped saving floor of Section~\ref{sec:assetdep}. The proof needs only that
$\Lambda_k$ is increasing in assets, which the clamp preserves.

\begin{proposition}[The clamped floor]\label{prop:clampfloor}
$\Lambda_k(z,a)\le A_k(a,z)$ on any reachable family, so capital per head is at least
$\sum_z\nu_z\Lambda_K(z,0)/(K+1)$, which is at least what Proposition~\ref{prop:affine} gives.
\formal{IncomeFluctuation.cohortFloorAt\_le\_cohortAssets},
\lean{IncomeFluctuation.le\_olgCapital\_at}
\end{proposition}

\subsection{A ceiling on capital}\label{sec:affceil}

Existence also needs capital supply below demand at the bottom of the interval, and there three
bounds are available, each asking for more and giving more.

The cheapest uses the budget and nothing else. A household cannot carry forward more than its cash
on hand, so an asset level $B$ with $RB+y_{\max}\le B$ is never left once entered.

\begin{proposition}[A forward-invariant level caps supply]\label{prop:ceiling}
If $0\le B\le\bar a$ and $RB+y_{\max}\le B$, then capital per head is at most $B$.
\formal{IncomeFluctuation.olgCapital\_le\_of\_invariant}
\end{proposition}

The second says the household will not want to carry that much. Consuming at least
$\kappa_k(m+H_k(z))$ by the constraint-incidence bound of Theorem~\ref{thm:incidence}, it saves at
most $(1-\kappa_k)(y_z+Ra)-\kappa_kH_k(z)$.

\begin{proposition}[The incidence ceiling]\label{prop:ceiling2}
If $B\le\bar a$ and, at every age and earnings state,
$(1-\kappa_j)(y_w+RB)-\kappa_jH_j(w)\le B$, then capital per head is at most $B$.
\formal{IncomeFluctuation.olgCapital\_le\_of\_incidence}
\end{proposition}

Both bound \emph{every} age's assets by one level, and so pay $(K+1)B$ for a cohort that holds
nothing when it is born, nothing when it dies, and anything like $B$ only in between. The floor
never made that mistake, and the third ceiling does not either: it is the same affine object,
built from the same slope $G$, with the saving ceiling in place of the saving floor,
\[
  \bar F_0=0,\qquad
  \bar F_{k+1}(z)=\textstyle\sum_{z'}\pi(z,z')\bar F_k(z')
    +G_k\bigl[(1-\kappa_{k+1})y_z-\kappa_{k+1}H_{k+1}(z)\bigr],
\]
with $H$ here the capped, risk-adjusted human wealth of Theorem~\ref{thm:incidence}.

\begin{proposition}[The affine ceiling]\label{prop:affceil}
$A_k(a,z)\le\bar F_k(z)+G_ka$, and hence capital per head is at most
$\sum_z\nu_z\bar F_K(z)/(K+1)$.
\formal{IncomeFluctuation.cohortAssets\_le\_cohortCeil},
\lean{IncomeFluctuation.olgCapital\_le\_cohortCeil}
\end{proposition}

Replacing a worst case over ages by an average over them is worth a factor of fourteen, and it
keeps working above the rate of time preference, where no forward-invariant level exists at all.

\begingroup\small
\begin{center}
\begin{tabular}{@{}lccccc@{}}
\toprule
& $r=-7\%$ & $r=-4\%$ & $r=0\%$ & $r=\lambda$ & $r=5\%$\\
\midrule
Capital demand & $56.25$ & $14.06$ & $7.03$ & $4.62$ & $4.33$\\
Proposition~\ref{prop:ceiling}, the budget & $42.52$ & $74.41$ & --- & --- & ---\\
Proposition~\ref{prop:ceiling2}, incidence & $22.17$ & $35.11$ & --- & --- & ---\\
Proposition~\ref{prop:affceil}, affine & $1.54$ & $2.28$ & $3.86$ & $5.98$ & $6.47$\\
True capital supply & $0.64$ & $1.03$ & $1.98$ & $3.78$ & $4.25$\\
\bottomrule
\end{tabular}
\end{center}
\endgroup

\subsection{Existence from primitives}

The two ends can now be put together, and the statement has taken four forms as the bounds
improved. All four are in the development; the last is the one to read.

\begin{theorem}[Existence from primitives]\label{thm:exist2}
Fix a rate interval and continuous demand $D$. Suppose there is a $B$ with $0\le B\le\bar a$,
$(1+r_{lo})B+y_{\max}\le B$ and $B\le D(r_{lo})$; and suppose $\rho_{K+1}<\bar a$ and
$D(r_{hi})\le\sum_z\nu_zF_K(z)/(K+1)$ at $r_{hi}$. Then there is a rate in the interval at which
the capital the cohorts accumulate equals the capital the firm demands.
\formal{IncomeFluctuation.exists\_olgEquilibrium\_rate},
\lean{IncomeFluctuation.exists\_olgEquilibrium\_of\_bounds}
\end{theorem}

Nothing in those hypotheses refers to the household's decisions. They are inequalities among the
discount factor, the interest rate, the earnings process, the horizon and the demand curve, with
$\rho$, $F$ and $G$ explicit finite recursions in those primitives. Three refinements follow.
Proposition~\ref{prop:slope} turns the floor into a closed form with no recursion left to run
(\lean{IncomeFluctuation.exists\_olgEquilibrium\_of\_geomBounds}, Theorem~\ref{thm:exist3} below);
and Proposition~\ref{prop:ceiling2} replaces the budget ceiling
(\lean{IncomeFluctuation.exists\_olgEquilibrium\_of\_sharpBounds}). The affine ceiling replaces it
again, and at that point no auxiliary $B$ and no invariance condition survive.

\begin{theorem}[Existence, both ends read the same way]\label{thm:exist4}
With the ceiling at $r_{lo}$ and the floor at $r_{hi}$ both taken as weighted averages of affine
bounds on a cohort's accumulated assets, a stationary equilibrium exists.
\formal{IncomeFluctuation.exists\_olgEquilibrium\_affine}
\end{theorem}

Reading the floor at the assets held as well gives the tightest form.

\begin{theorem}[Existence with both ends read at the assets held]\label{thm:exist5}
With the affine ceiling at $r_{lo}$ and the clamped floor at $r_{hi}$, a stationary equilibrium
exists.
\formal{IncomeFluctuation.exists\_olgEquilibrium\_clamped}
\end{theorem}

The two boundary conditions are then the same object built twice, once from the saving ceiling and
once from the saving floor, sharing a slope. Each half asks for what it can have: at the bottom
the artefactual cap is slack, which holds below the rate of time preference; at the top it cannot
be, and the reachable family of Section~\ref{sec:reach} stands in. So the low endpoint must sit
below $\lambda$ and the high endpoint above it, and Section~\ref{sec:sharpnum} shows a bracket that
does.

For completeness, the closed-form variant.

\begin{theorem}[Existence in closed form]\label{thm:exist3}
Under the conditions of Theorem~\ref{thm:exist2}, with the floor hypothesis replaced by the
closed-form expression of Proposition~\ref{prop:slope} and with $1\le\Thorn\le R$ at $r_{hi}$, a
stationary equilibrium exists.
\formal{IncomeFluctuation.exists\_olgEquilibrium\_of\_geomBounds}
\end{theorem}

\subsection{What the numbers say}\label{sec:num2}

The calibration is the one used throughout: $J=60$, $\beta=0.96$, relative risk aversion $3$, the
Tauchen discretisation of Aiyagari's earnings process with $\rho=0.9$ and $\sigma=0.4$, and a
Cobb--Douglas technology with $\alpha=9/25$ and $\delta=2/25$, so capital demand is
$\alpha/((1-\alpha)(r+\delta))$ and $\lambda=4.17\%$. Capital supply is the average asset holding
across the sixty cohorts, each born with nothing. The last row is not the solved supply but the
floor of Proposition~\ref{prop:affine} (\texttt{numerics/olgexist.m}).

\begingroup\small
\begin{center}
\begin{tabular}{@{}lccccccc@{}}
\toprule
$r$ & $-4\%$ & $0\%$ & $\lambda$ & $6\%$ & $8\%$ & $10\%$ & $15\%$ \\
\midrule
Capital demand & $14.06$ & $7.03$ & $4.62$ & $4.02$ & $3.52$ & $3.13$ & $2.45$ \\
Capital supply & $1.03$ & $1.98$ & $3.78$ & $4.85$ & $6.18$ & $7.69$ & $12.49$ \\
Provable floor & $0.00$ & $0.00$ & $0.00$ & $1.52$ & $3.13$ & $4.75$ & $9.40$ \\
\bottomrule
\end{tabular}
\end{center}
\endgroup

Supply crosses demand at $5.09\%$, above the rate of time preference, which is what
Section~\ref{sec:reach} has to accommodate; and the floor is flat below $\lambda$ and climbs above
it, which is Theorem~\ref{thm:patient}. The optimal path reaches assets of $82$ at $r=10\%$,
against the crude reachable bound of $9{,}033$; the constants are ordinary, unlike the stationary
case, where the corresponding instantiation put $r_{hi}$ within $10^{-27}$ of $\lambda$.

\subsection{The theorem checked, hypothesis by hypothesis}\label{sec:sharpnum}

Take $r_{lo}=2\%$, $r_{hi}=8\%$ and an artefactual cap of $10^4$. Every hypothesis of
Theorem~\ref{thm:exist5} holds (\texttt{numerics/clampcheck.m}): the cap is forward invariant at
$r_{lo}$ with margin $134$; the ceiling is $4.81$ against a demand of $5.63$; a lifetime of
feasible saving reaches $3{,}730$ against the cap; and the floor is $3.80$ against a demand of
$3.52$. The bracket contains the true equilibrium at $5.09\%$.

Sweeping the endpoints shows what each bound is worth. The low endpoint may be taken as high as

\begingroup\small
\begin{center}
\begin{tabular}{@{}lc@{\qquad}lc@{}}
\toprule
Ceiling used & highest $r_{lo}$ & Floor used & lowest $r_{hi}$\\
\midrule
Proposition~\ref{prop:ceiling}, the budget & $-6.73\%$
  & closed form, crude $G_k\ge1$ & $610\%$\\
Proposition~\ref{prop:ceiling2}, incidence & $-5.85\%$
  & closed form, Proposition~\ref{prop:slope} & $18.2\%$\\
Proposition~\ref{prop:affceil}, affine & $+2.78\%$
  & Proposition~\ref{prop:affine}, affine & $8.38\%$\\
& & Proposition~\ref{prop:clampfloor}, clamped & $7.67\%$\\
\bottomrule
\end{tabular}
\end{center}
\endgroup

So the tightest bracket the development gives is $[2.78\%,\,7.67\%]$, and the absolute limit is
$5.09\%$ from both sides, since no bound on capital supply can cross demand before supply does.
The closed-form floor is reported because it needs no recursion at all; at a calibrated rate it is
much the weaker, giving $0.48$ at $r=10\%$ against a true floor of $4.75$.

Neither bracket is tight, and the same fact stops them at both ends. Appendix~\ref{app:routes} gives
it: the obstruction at the top and the obstruction at the bottom are one obstruction, that a power
mean with negative exponent is superadditive rather than $1$-Lipschitz.
